%%%%%%%%%%%%%%%%%%%%%%%%%%%%%%%%%%%%%%%%%%%%%%%%%%%%%%%%%%%%%%%%%%%%%%%%%%%%%%%%
%%
%% Abstract of your thesis in default document language (English)
%%
%%%%%%%%%%%%%%%%%%%%%%%%%%%%%%%%%%%%%%%%%%%%%%%%%%%%%%%%%%%%%%%%%%%%%%%%%%%%%%%%

\cleardoubleoddpage%  Make sure to start abstract on a new odd page
\begin{abstract*}
Audio deepfake detectors are routinely reported as effective on
standard benchmark corpora, yet their behaviour under a
simultaneous shift in speaker dialect, recording domain, and
synthesis generator is much less well characterised. This thesis
studies how detector performance changes across two English
corpora that differ in dialect and recording style, DECTE
(Diachronic Electronic Corpus of Tyneside English) and VCTK
(Voice Cloning Toolkit corpus), and across two neural speech
generators, XTTS v2 and OpenVoice v2. Two detectors are used: an
AASIST-style graph-attention system (the AuralGuard baseline) and
a lightweight LFCC + logistic-regression system as a
cross-architecture check. Detector behaviour is summarised with
equal error rate (EER), AUC, and bootstrap 95\,\% confidence
intervals.

The two generators do not agree. On XTTS v2 spoofs, both
detectors show a positive DECTE-minus-VCTK EER gap:
$+12.80$\,pp on AASIST and $+12.44$\,pp on LFCC + LR. On
OpenVoice v2 spoofs the direction reverses on both detectors:
$-26.74$\,pp on AASIST and $-42.62$\,pp on LFCC + LR. The pattern
is therefore a generator-specific corpus interaction: a
dialect / domain gap whose sign depends on which generator
produced the spoof, not a pure dialect bias intrinsic to the
detector.

A partial-unfreeze controlled adaptation of the AuralGuard
backbone (mitigation v2) reduced DECTE XTTS EER from 40.70\,\%
to 23.26\,\%, a paired delta of $-17.44$\,pp with bootstrap
95\,\% CI $[-28.49,\,-9.30]$, and produced no clear VCTK
regression on either generator. Subgroup diagnostics on the
mitigation-held-out DECTE slice show broad but uneven
per-subgroup gains and are reported as descriptive diagnostics
only, not a fairness audit. Robust evaluation of audio deepfake
detectors must therefore vary corpus, generator, and detector
architecture jointly.
\end{abstract*}

%%
%%%%%%%%%%%%%%%%%%%%%%%%%%%%%%%%%%%%%%%%%%%%%%%%%%%%%%%%%%%%%%%%%%%%%%%%%%%%%%%%


%%%%%%%%%%%%%%%%%%%%%%%%%%%%%%%%%%%%%%%%%%%%%%%%%%%%%%%%%%%%%%%%%%%%%%%%%%%%%%%%
%%
%% Abstract of your thesis in an additional language (German / Kurzfassung)
%%
%%%%%%%%%%%%%%%%%%%%%%%%%%%%%%%%%%%%%%%%%%%%%%%%%%%%%%%%%%%%%%%%%%%%%%%%%%%%%%%%

\cleardoubleoddpage%  Make sure to start abstract on a new odd page
\begin{otherlanguage}{ngerman}
\begin{abstract*}
Audio-Deepfake-Detektoren werden auf g\"angigen Benchmark-Korpora
h\"aufig als leistungsf\"ahig berichtet; ihr Verhalten unter einer
gleichzeitigen Verschiebung von Dialekt, Aufnahmedom\"ane und
Synthesegenerator ist bislang jedoch deutlich weniger untersucht.
Diese Arbeit untersucht, wie sich die Detektor-Leistung \"uber
zwei englischsprachige Korpora hinweg ver\"andert, die sich in
Dialekt und Aufnahmesituation unterscheiden: DECTE
(Diachronic Electronic Corpus of Tyneside English) und VCTK
(Voice Cloning Toolkit corpus). Sie vergleicht die
Detektor-Leistung au{\ss}erdem \"uber zwei neuronale
Sprachgeneratoren, XTTS v2 und OpenVoice v2. Zwei Detektoren
werden verwendet: ein AASIST-artiges Graph-Attention-System
(die AuralGuard-Baseline) und ein leichtgewichtiges
LFCC + logistische-Regression-System als architektur\"ubergreifende
Kontrolle. Die Detektor-Leistung wird mit Equal Error Rate
(EER), AUC und Bootstrap-95\,\%-Konfidenzintervallen
zusammengefasst.

Die beiden Generatoren zeigen kein einheitliches Bild. Auf
XTTS-v2-Spoofs weisen beide Detektoren eine positive
DECTE-minus-VCTK-EER-Differenz auf: $+12{,}80$\,pp auf AASIST
und $+12{,}44$\,pp auf LFCC + LR. Auf OpenVoice-v2-Spoofs
kehrt sich die Richtung auf beiden Detektoren um:
$-26{,}74$\,pp auf AASIST und $-42{,}62$\,pp auf LFCC + LR.
Das Muster wird daher als generator-spezifische Korpus-Interaktion
beschrieben: als eine Dialekt-/Dom\"anen-L\"ucke, deren Vorzeichen
davon abh\"angt, welcher Generator den Spoof erzeugt hat. Es
handelt sich nicht um einen reinen, dem Detektor eigenen
Dialekt-Bias.

Eine partielle Unfreeze-Anpassung des AuralGuard-Backbones
(Mitigation v2) reduzierte den DECTE-XTTS-EER von
40{,}70\,\% auf 23{,}26\,\%, ein paarweises Delta von
$-17{,}44$\,pp mit Bootstrap-95\,\%-KI $[-28{,}49,\,-9{,}30]$,
ohne erkennbare VCTK-Regression auf beiden Generatoren.
Subgruppen-Diagnostiken auf dem f\"ur die Mitigation
zur\"uckgehaltenen DECTE-Slice zeigen breite, aber
ungleichm\"a{\ss}ige Gewinne und sind ausschlie{\ss}lich als
deskriptive Diagnostik, nicht als Fairness-Audit zu verstehen.
Robuste Bewertung von Audio-Deepfake-Detektoren muss
demnach Korpus, Generator und Detektor-Architektur gemeinsam
variieren.
\end{abstract*}
\end{otherlanguage}

%%
%%%%%%%%%%%%%%%%%%%%%%%%%%%%%%%%%%%%%%%%%%%%%%%%%%%%%%%%%%%%%%%%%%%%%%%%%%%%%%%%
